\documentclass[11pt,a4paper]{report}
%\renewcommand{\rmdefault}{phv} % Arial
%\renewcommand{\sfdefault}{phv} % Arial
\usepackage{helvet}
\renewcommand{\familydefault}{\sfdefault}
\setlength{\parskip}{1em}
\usepackage{graphicx}
\usepackage{titlepic}
\usepackage{index}
\makeindex
\usepackage{listings}
\usepackage{color}

\definecolor{dkgreen}{rgb}{0,0.6,0}
\definecolor{gray}{rgb}{0.5,0.5,0.5}
\definecolor{mauve}{rgb}{0.58,0,0.82}

\lstset{frame=tb,
  aboveskip=3mm,
  belowskip=3mm,
  showstringspaces=false,
  columns=flexible,
  basicstyle={\small\ttfamily},
  numbers=none,
  numberstyle=\tiny\color{gray},
  keywordstyle=\color{blue},
  commentstyle=\color{dkgreen},
  stringstyle=\color{mauve},
  breaklines=true,
  breakatwhitespace=true,
  tabsize=3
}

\title{\includegraphics[scale=1.5]{twt}\\~\\~\\ \LARGE{\textbf{TWT Programming
Language}}\\~\\ \large{(Talagang Walang Tulog)}\\~\\ \large{Language Design and Functionalities}}
\date{~\\ December 2015}
\author{Jahziel Rae Arceo\\ Gabriel Kelly Navarro\\ Mikaela Jun Lenon\\ Maria
Rosario Gueco\\ Marbille Juntado\\~\\~\\ CS 150 HTWX}

\begin{document}

\maketitle

\tableofcontents

\chapter{Introduction}

\section{The TWT programming language}

\index*{TWT} is simply a programming language designed out of curiosity. There
were really many challenges of deciding on which fits best on the category of a
usable language, which the designers took so much time for considering and
creating such.

\subsection{Language Name}

Halfway before the appropriate parser was created for the
language, there was no appropriate meaning that the programmers can imply on
the initials TWT, except that it was Twitter-inspired.
However, since they were inspired by the countless nights that they weren't of
good sleep, they made a way for the meaning of the name. TWT refers to "Talagang Walang Tulog",
or "really no sleep", the phenomenon that happened to them during the creation
of the language.

\subsection{Paradigm}

The TWT programming language is imperative and functional. The programming
language works in a fashion nearly similar to C, Python, and Pascal. Function
declarations are C-styled. Also, keywords that begin and end a program are used
in the language, which seems to be like Pascal. The language also borrows features from Python, which was the
language used to create the new language.

\subsection{Inspiration}

\index*{Twitter}, a social networking site, was the main inspiration on creating
the language. As seen on the next chapter, the keywords and reserved words of the
language mostly came from the jargons seen on the website. \par

Other than Twitter, the programmers also made the intricacy of programming an
inspiration to create the language. Syntaxes on the language are simpler
compared to other languages, which makes new programmers learn the language
easily.

\chapter{Grammar Definition}

In this chapter, we get to learn about the grammar of the language using its
Backus-Naur form (BNF). By this, we can get a hint of what the language looks
like on encoding.

\section{Identifiers}

The following conventions apply to \index*{variables} and values entered in the
programming language:

\begin{lstlisting}
<INT> ->				<Digit><INT> | <Digit>
<Digit> ->			"0" | "1" | "2" | "3" | "4" | "5" | "6"
						| "7" | "8" | "9"

<CHAR> ->			"'"<ASCII>"'"

<FLOAT> ->			<Digit>"."<Digit>
						(Note: Strictly followed)

<STRING> -> 		"""<STRINGEXT>"""

<STRINGEXT> -> 	<CHAR><STRINGEXT> | <CHAR> | EPSILON

<VARIABLE> -> 		<ALPHABET><VARIABLEEXT>

<VARIABLEEXT> -> <ALPHANUMERIC><VARIABLEEXT>
						| <ALPHANUMERIC> | EPSILON
\end{lstlisting}

Note that the grammar for float type values is strict in the programming
language, so it will not accept an integer as float. The data must explicitly
indicate that it is indeed a floating-point value.

\section{Data Types}

The TWT language is statically-typed. The programmer should explicitly state the
data type he is to use.

Data types were defined in the language grammar as follows:

\begin{lstlisting}
<Dtype> -> "@INT" | "@CHIRP" | "@COKE" | "@MSG" | "@TRALSE"
\end{lstlisting}

Data types are as follows:
\begin{itemize}
  \item \textbf{@INT:} integer
  \item \textbf{@CHIRP:} character (ASCII)
  \item \textbf{@COKE:} float
  \item \textbf{@MSG:} string
  \item \textbf{@TRALSE:} boolean
\end{itemize}

The '@' symbol is used to easily identify words as types. By inspiration, it
takes the @-mention paradigm of Twitter as its data type recognition
mechanism.

\section{Expressions and Assignment Statements}

The following rules apply to \index*{expressions} and \index*{assignments}:
\begin{lstlisting}
<Assignment> -> 	<DType> VARIABLE "=" VARIABLE
                	| <DType> VARIABLE "=" <Reading>
                	| "@INT" VARIABLE "=" INT
                	| "@CHIRP" VARIABLE "=" CHAR
                	| "@COKE" VARIABLE "=" FLOAT
                	| "@MSG" VARIABLE "=" STRING
                	| "@TRALSE" VARIABLE "=" "YES"
                	| "@TRALSE" VARIABLE "=" "NO"
                	| <DType> VARIABLE "=" "(" <Exp> ")"

<Reading> -> 		"REPLY" VARIABLE

<Exp> -> 			<Term> <ExpPrime>

<ExpPrime> ->		"+" <Term> <EspPrime>
            		| "-" <Term> <ExpPrime>
            
<Term> -> 			<Fact> <TermPrime>

<TermPrime> ->		"*" <Fact> <TermPrime>
            		| "/" <Fact> <TermPrime>
            
<Fact> -> 			INT | CHAR | FLOAT | VARIABLE | "(" <Exp> ")"

\end{lstlisting}

\section{Statement Level Control Structures}

The following rules apply to control structures available in the language:
\begin{lstlisting}
<State> ->			<Loop> "#" <StatePrime>
        				| <If> "#" <StatePrime>
        				| <Assignment> "#" <StatePrime>
        				| <Call> "#" <StatePrime>
        				| <Printing> "#" <StatePrime>
        				| <Control> "#" <StatePrime>
        			
<StatePrime> -> 	<Loop> "#" <StatePrime>
        				| <If> "#" <StatePrime>
        				| <Assignment> "#" <StatePrime>
        				| <Call> "#" <StatePrime>
        				| <Printing> "#" <StatePrime>
        				| <Read> "#" <StatePrime>
        				| <Control> "#" <StatePrime>
        			
<Loop> -> 			"RT" <Condition> {" <Block> "}"

<If> -> 				"IF" <Condition> "FOLLOW" "{" <Block> "}"
						<ElseIf>
						<Else>
        				| "IF" <Condition> "FOLLOW" "{" <Block> "}"
        				<Else>
        				| "IF" <Condition> "FOLLOW" "{" <Block> "}"
        			
<ElseIf> ->			"ELSEIF" <Condition> "FOLLOW" "{" <Block> "}"
						<Elseif>

<Else> -> 			"ELSE" <Condition> "FOLLOW" "{" <Block> "}"
\end{lstlisting}

\section{Subprograms}

This is how functions or subprograms are declared in the language:
\begin{lstlisting}
<Declaration> ->	<Dtype> VARIABLE "(" <Args> ")"
						"{" <Block> <Return> "}"
						| <Dtype> VARIABLE "(" <Args> ")" "{" <Block> "}"
\end{lstlisting}

\section{Other syntaxes not listed on previous sections}

\begin{lstlisting}
<Program> ->		<Declaration> <Main>
            		| <Main>
            
<Main> ->			"LOGIN" <Block> "LOGOUT"
          			| "LOGIN" "LOGOUT"
          			
<Block> ->			<State>

<Call> ->			"HOOT "VARIABLE "(" <Args> ")"

<Printing> -> 		"TWEET" VARIABLE
              		| "TWEET" "(" <EXP> ")"
              		| "TWEET" INT
              		| "TWEET" FLOAT
              		| "TWEET" CHAR
              		| "TWEET" STRING
              		| "TWEET" TRUE
              		| "TWEET" FALSE
<Reading> ->		"REPLY" VARIABLE

<Control> ->		"UNFOLLOW" | "LIKE" | "BLOCK"

<Condition> ->		"(" <Conditional> ")" | "~" <Condition>

<Conditional> ->	VARIABLE <Conditional> VARIABLE
          			| { VARIABLE | INT | CHAR | FLOAT } <CondOp> <Condition>
          			| <Conditional> <CondOp> <Conditional>
          			| <Condition> <CondOp> { VARIABLE | INT | CHAR | FLOAT }
          			| { VARIABLE | INT | CHAR | FLOAT } <CondOp>
          			{ VARIABLE | INT | CHAR | FLOAT }
          			
<CondOp> ->			">=" | "<=" | "==" | ">" | "<"

<Args> ->			<Dtype> <Vname> "," <Args>

<Return> ->			"REPORT" { INT | CHAR | FLOAT | STRING | VARIABLE }
						| "REPORT" <Exp>
\end{lstlisting}

\chapter{Lexical and Syntax Analysis}

\section{Parser}

The \index*{parser} for the language uses LR parsing, so no parse tables were
used to parse the strings for the language. By that, it is obvious in the grammar that
revisions for left recursion were done, so that no conflict will happen on the
parsing process. \par

On the next lines of code is the Python code for the lexical
analyzer\index{lexical analyzer}, together with the token numbers
assigned to them.

\begin{lstlisting}
EOF = -1
int_dec = 3
float_dec = 4
char_dec =  5
string_dec = 6
bool_dec = 7
login = 1
logout = 0
if_state = 20
elseif_state = 21
else_state = 22
loop_state = 30
break_state = 31
continue_state = 32
exit_state = 33
exec_state = 34
read_state = 10
print_state = 11
return_state = 12
TRUE = 13
FALSE = 14
NOT = 15
INT = 71
FLOAT = 72
CHAR = 73
STRING = 74
openBrace = 41
closeBrace = 42
plusSign = 43
minusSign = 44
divSign = 45
mulSign = 46
asSign = 47
openParen = 48
closeParen = 49
ENDOFSTATE = 50
commaSign = 51
lessEqSign = 52
greatEqSign = 53
eqSign = 54
lesserSign = 55
greaterSign = 56
VARIABLE = 100

nextToken = 0

def isfloat(str):
    try:
        float(str)
    except ValueError:
        return False
    return True

def lex():
    global input
    global nextToken
    nextString = input[0]
    if nextString == 'EOF':
        nextToken = EOF
    elif nextString == 'LOGIN':
        nextToken = login
    elif nextString == 'LOGOUT':
        nextToken = logout
    elif nextString == '@INT':
        nextToken = int_dec
    elif nextString == '@COKE':
        nextToken = float_dec
    elif nextString == '@CHIRP':
        nextToken = char_dec
    elif nextString == '@MSG':
        nextToken = string_dec
    elif nextString == '@TRALSE':
        nextToken = bool_dec
    elif nextString == 'IF':
        outfile.write('if ')
        nextToken = if_state
    elif nextString == 'ELSEIF':
        outfile.write('elif ')
        nextToken = elseif_state
    elif nextString == 'ELSE':
        outfile.write('else ')
        nextToken = else_state
    elif nextString == 'RT':
        outfile.write('while ')
        nextToken = loop_state
    elif nextString == 'UNFOLLOW':
        outfile.write('break')
        nextToken = break_state
    elif nextString == 'LIKE':
        outfile.write('continue')
        nextToken = continue_state
    elif nextString == 'BLOCK':
        outfile.write('exit')
        nextToken = exit_state
    elif nextString == 'FOLLOW':
        nextToken = exec_state
    elif nextString == 'REPLY':
        nextToken = read_state
    elif nextString == 'TWEET':
        outfile.write('print ')
        nextToken = print_state
    elif nextString == 'REPORT':
        nextToken = return_state
    elif nextString == 'YES':
        nextToken = TRUE
    elif nextString == 'NO':
        nextToken = FALSE
    elif nextString == '~':
        nextToken = NOT
    elif nextString.isdigit():
        outfile.write(nextString)
        nextToken = INT
    elif nextString[0] == '-' and nextString[1:].isdigit():
        outfile.write(nextString)
        nextToken = INT
    elif isfloat(nextString):
        outfile.write(nextString)
        nextToken = FLOAT
    elif (nextString[0] == '\'') and (nextString[-1] == '\'') and (nextString.len() == 1):
        outfile.write(nextString)
        nextToken = CHAR
    elif (nextString[0] == '\"') and (nextString[-1] == '\"'):
        outfile.write(nextString)
        nextToken = STRING
    elif nextString == '{':
        nextToken = openBrace
    elif nextString == '}':
        nextToken = closeBrace
    elif nextString == '+':
        outfile.write(nextString)
        nextToken = plusSign
    elif nextString == '-':
        outfile.write(nextString)
        nextToken = minusSign
    elif nextString == '/':
        outfile.write(nextString)
        nextToken = divSign
    elif nextString == '*':
        outfile.write(nextString)
        nextToken = mulSign
    elif nextString == '=':
        outfile.write(nextString)
        nextToken = asSign
    elif nextString == '(':
        outfile.write(nextString)
        nextToken = openParen
    elif nextString == ')':
        outfile.write(nextString)
        nextToken = closeParen
    elif nextString == '#':
        outfile.write('\n')
        nextToken = ENDOFSTATE
    elif nextString == ',':
        outfile.write(nextString)
        nextToken = commaSign
    elif nextString == '>=':
        outfile.write(nextString)
        nextToken = lessEqSign
    elif nextString == '<=':
        outfile.write(nextString)
        nextToken = greatEqSign
    elif nextString == '==':
        outfile.write(nextString)
        nextToken = eqSign
    elif nextString == '>':
        outfile.write(nextString)
        nextToken = lesserSign
    elif nextString == '<':
        outfile.write(nextString)
        nextToken = greaterSign
    else:
        outfile.write(nextString)
        nextToken = VARIABLE
    print("Next token is: " + str(nextToken))
    print("Next string is: " + nextString)
    del input[0]

\end{lstlisting} 

All reserved words are explicitly stated in the analyzer.

\chapter{Names, Binding, and Scoping}

\section{Case sensitivity}

The TWT programming language was designed in a way that the user must
exactly write the language in whichever case it is. By that, TWT is
case-sensitive. For example, \textbf{@int} is not the same with \textbf{@INT}.

\section{Reserved words}

The following is a list of the reserved words in the language:

\begin{center}
\begin{tabular}{ c c c c c }
 LOGIN & LOGOUT & REPORT & TWEET & IF \\ 
 ELSEIF & ELSE & HOOT & REPLY & UNFOLLOW \\  
 @INT & @CHIRP & @COKE & @MSG & @TRALSE \\
 LIKE & BLOCK & YES & NO
\end{tabular}
\end{center}

\section{Name form}

By convention, as explained in the grammar, the \index*{names} must be
alphabetic at the start of each string, for variables. The language follows C-style rules.

\section{Binding}

The language follows Python-based \index*{binding}, but on the programming
language scope, it was designed to be statically bound. As noticed, Python does dynamic
binding, however, the programmers chose to require having types specified before
each run.

\section{Lifetime and scope}

Conventions were also followed on variable lifetime and scope. That is, what
Python generally has on variable lifetime and scope applies to TWT as well.

\section{Blocks}

In TWT, blocks can be determined by curly braces, as what C has. Note the
following code example:

\begin{lstlisting}
RT ( X < 4 ) { 
		X = ( X + 1 ) #
		Y = ( Y + 2 ) #
		TWEET X #
 } #
\end{lstlisting}

The example above is a loop. It has the three statements "X = ( X + 1 ) \#", "Y
= ( Y + 2 ) \#", and "TWEET X \#" in the block example.

\chapter{Data types}

\section{Primitive data types}

As referenced to the grammar in Chapter 2, these are the data types\index{data
types} available in the language:

\begin{itemize}
  \item \textbf{@INT:} integer
  \item \textbf{@CHIRP:} character (ASCII)
  \item \textbf{@COKE:} float
  \item \textbf{@MSG:} string
  \item \textbf{@TRALSE:} boolean
\end{itemize}

\section{Strings}

The programming language supports \index*{strings}, and is a primitive data
type.
Like C, the strings are determined by double quotation marks.\par

There are no limitations on string declarations, and everything that applies to
every other language also applies to TWT.

\section{Note on user-defined types and type-checking}

As of writing, \index*{user-defined typing} is not available in the language.
The user must strictly use the available data types.\par

On type-checking and coercion, Python conventions are used, that is, the
features that Python has are used in the language.

\chapter{Expression and Assignment Statements}

\section{Arithmetic expressions and overloading}

This is the list of \index{mathematical symbols}mathematical symbols used in the
language.

\begin{center}
\begin{tabular}{ c c c c }
 + & - & * & / \\
 = & ( & )
\end{tabular}
\end{center}

\noindent Listed are its functions:

\begin{center}
\begin{tabular}{ c l }
  \textbf{+} & addition\\
  \textbf{-} & subtraction\\
  \textbf{*} & multiplication\\
  \textbf{/} & division\\
  \textbf{=} & assignment\\
  \textbf{(} & grouping (opening)\\
  \textbf{)} & grouping (closing)
\end{tabular}
\end{center}

Also the language is using the infix notation for expressions. For the
operators, there is no overloading. Operations are defined as to what they
represent.

\section{Boolean expressions}

This is the list of \index{boolean symbols}boolean symbols used in the language.

\begin{center}
\begin{tabular}{ c c c c }
 YES & NO & == & \textless = \\
 \textgreater = & \textless & \textgreater & \textasciitilde
\end{tabular}
\end{center}

\noindent Listed are its functions:

\begin{center}
\begin{tabular}{ c l }
  \textbf{YES} & TRUE\\
  \textbf{NO} & FALSE\\
  \textbf{==} & equal\\
  \textbf{\textless}= & less than or equal\\
  \textbf{\textgreater}= & greater than or equal\\
  \textbf{\textless} & less than\\
  \textbf{\textgreater} & greater than\\
  \textbf{\textasciitilde} & NOT
\end{tabular}
\end{center}

As what applies in most languages on boolean statements, the conventional
approach to boolean expressions apply to TWT, specifically the operator
precedence.

\chapter{Statement-Level Control Structures}

On the sections to this chapter, these lines of code will be referred to.

\begin{lstlisting}
RT ( X < 4 ) {
		X = ( X + 1 ) #
		Y = ( Y + 2 ) #
		TWEET X #
    IF ( X >= Y ) FOLLOW {
      TWEET "Ahihi" #
    }
    ELSE {
      Y = X #
      HOOT HAHAHA ( @INT X ) #
    } #
    RT ( X < 4 ) {
    		X = ( X + 1 ) #
    		Y = ( Y + 2 ) #
    		TWEET X #
        IF ( X >= Y ) FOLLOW {
          TWEET "Ahihi" #
        }
        ELSE {
          Y = X #
          HOOT HAHAHA ( @INT X ) #
        } #
     } #
 } #
\end{lstlisting}

\section{Selection statements}

As seen on the code before this section, IF and ELSE blocks are available as
selection statements\index{selection statements}. Also some features, like the
ELSEIF block, is also available.

For ELSEIF usage, refer to this code (an edited version of the code at the
beginning of the chapter):

\begin{lstlisting}
IF ( X >= Y ) FOLLOW {
  TWEET "Ahihi" #
}
ELSEIF ( X <= Y - 5 ) FOLLOW {
  TWEET "Huhuhu" #
}
ELSE {
  Y = X #
  HOOT HAHAHA ( @INT X ) #
} #
\end{lstlisting}

\section{Iterative statements}

Only one kind of iterative statement is available in TWT, which is RT. As seen
on the beginning of this chapter, the statement can have its condition. This
feature is equivalent to the while loop of other languages.

\chapter{Subprograms}

In reference to Chapter 2, this is how \index*{subprograms} are declared in TWT.

\begin{lstlisting}
<Dtype> NAME ( <Args> ) {
	<Block>
	<Return>
}

or

<Dtype> NAME ( <Args> ) {
	<Block>
}
\end{lstlisting}

The first variation recognizes a return statement at the end of the function.
Note that it must be on the last part of the block. This is considered good
programming practice.\par

The second variation is the same version but without the return statement. In
essence, you can create a subprogram without return values. Note too that even
if \textless Dtype\textgreater\ is on the grammar rules, the only data type
associated with subprograms which have no return value are the void subprogram return types.

\section{Parameter passing and function calling}

The keyword HOOT is used to call subprograms. A code sample is shown on the next
page to see how it works.

\begin{lstlisting}
@INT HAHAHA ( @INT X ) {
    TWEET "Kaya mo yan" #
    TWEET "Oo Kaya mo yan" #
    @INT X = 0 #
}

@INT HIHIHI ( @INT Y ) {
    TWEET "Kaya mo yan" #
    TWEET "Oo Kaya mo yan" #
    @INT Y = 0 #
}

/...statements here.../
HOOT HAHAHA ( @INT X ) #
HOOT HIHIHI ( @INT X ) #
/...statements here.../
\end{lstlisting}

As seen on the code example, HOOT is entered before a function name. That is to
signal the interpreter that a function call is being done. Also, parameters are
inserted inside the parentheses. They can be separated by commas if they are
multiple.

\section{Recursion}

Recursion can be also executed through HOOT. A HOOT call inside a function
commences \index*{recursion}.

\printindex

\end{document}
